
\documentclass{article}
\usepackage{colortbl}
\usepackage{makecell}
\usepackage{multirow}
\usepackage{supertabular}

\begin{document}

\newcounter{utterance}

\centering \large Interaction Transcript for game `dond', experiment `coop\_en', episode 14 with Qwen3{-}32B{-}t0.0.
\vspace{24pt}

{ \footnotesize  \setcounter{utterance}{1}
\setlength{\tabcolsep}{0pt}
\begin{supertabular}{c@{$\;$}|p{.15\linewidth}@{}p{.15\linewidth}p{.15\linewidth}p{.15\linewidth}p{.15\linewidth}p{.15\linewidth}}
    \# & $\;$A & \multicolumn{4}{c}{Game Master} & $\;\:$B\\
    \hline

    \theutterance \stepcounter{utterance}  
    & & \multicolumn{4}{p{0.6\linewidth}}{
        \cellcolor[rgb]{0.9,0.9,0.9}{
            \makecell[{{p{\linewidth}}}]{
                \texttt{\tiny{[P1$\langle$GM]}}
                \texttt{You are playing a cooperative negotiation game in which you have to agree on how to divide a set of items among you and another player.} \\
\\ 
\texttt{Rules:} \\
\texttt{(a) You and the other player are given a set of items. Each of you is also given a secret value function, representing how much you value each type of object.} \\
\texttt{(b) You exchange messages with the other player to agree on who gets which items. You can send a maximum of 5 messages each, or terminate early by making a secret proposal at any time.} \\
\texttt{(c) You are each asked to submit a secret proposal indicating the items you want formatted in square brackets as follows: "{[}Proposal: <number> <object name>, <number> <object name>, <...>{]}"} \\
\texttt{(d) If your proposals are complementary, i.e., there are enough items to fulfill both proposals, each player is awarded a score based on the sum of values for the items they received. Otherwise, both of you get zero points.} \\
\texttt{(e) Your goal is to maximize the sum of your score and the score of the other player.} \\
\\ 
\texttt{Let us start.} \\
\\ 
\texttt{The set of available items is:} \\
\\ 
\texttt{2 mirrors, 1 brick, 3 knives, 1 comb} \\
\\ 
\texttt{Your secret value function is:} \\
\\ 
\texttt{mirror: 1, brick: 4, knife: 1, comb: 1} \\
\\ 
\texttt{IMPORTANT: Your messages, unless it is the secret proposal, are directly transmitted to the other player, so do not include any response to the rules or text announcing your message. To make a secret proposal, use the indicated format. Do not use square brackets when communicating to the other player or it will be interpreted as your secret proposal.} \\
            }
        }
    }
    & \\ \\

    \theutterance \stepcounter{utterance}  
    & \multicolumn{4}{p{0.6\linewidth}}{
        \cellcolor[rgb]{0.9,0.9,0.9}{
            \makecell[{{p{\linewidth}}}]{
                \texttt{\tiny{[P1$\rangle$GM]}}
                \texttt{<think>} \\
\texttt{Okay, let's see. I need to divide these items between me and another player. The items are 2 mirrors, 1 brick, 3 knives, 1 comb. My value function is mirror:1, brick:4, knife:1, comb:1. The other player's values are secret, so I don't know what they value. My goal is to maximize the total score (mine plus theirs) by making sure our proposals are complementary.} \\
\\ 
\texttt{First, I should think about what items I value the most. The brick is worth 4 to me, which is the highest. Then mirrors, knives, and comb are all 1 each. So I probably want the brick. But if the other player also values the brick highly, maybe we need to split it. But since I don't know their values, I have to make an educated guess.} \\
\\ 
\texttt{The other player might have different priorities. For example, maybe they value knives more, or the comb. Since there are 3 knives, maybe splitting them could be beneficial. But I need to make sure that whatever I propose, there's enough items left for both of us.} \\
\\ 
\texttt{If I take the brick (4 points for me), then the other player might get the rest. But if they also want the brick, that could cause a conflict. Alternatively, maybe I should take the brick and some other items, but I need to leave enough for them. Let's see: if I take the brick, then there are 2 mirrors, 3 knives, 1 comb left. If the other player wants those, maybe they can take some of the knives and mirrors. Since I value knives at 1 each, maybe they value them higher. So if I take the brick and maybe one knife, then the other player can take the rest. But I need to make sure that the total is maximized.} \\
\\ 
\texttt{Alternatively, maybe the other player values the comb more. If I take the brick and leave the comb, that might be better. But I don't know. Since my main high value is the brick, I should probably secure that. Let me think about possible proposals.} \\
\\ 
\texttt{If I propose {[}Proposal: 1 brick, 1 mirror, 1 knife{]}, that would give me 4 +1 +1 = 6. Then the other player could take the remaining 1 mirror, 2 knives, 1 comb. If their values are such that they can get a good} \\
            }
        }
    }
    & & \\ \\

    \theutterance \stepcounter{utterance}  
    & & & \multicolumn{2}{p{0.3\linewidth}}{
        \cellcolor[rgb]{0.9,0.9,0.9}{
            \makecell[{{p{\linewidth}}}]{
                \texttt{\tiny{[GM$|$GM]}}
                \texttt{proposal} \\
            }
        }
    }
    & & \\ \\

    \theutterance \stepcounter{utterance}  
    & & & \multicolumn{2}{p{0.3\linewidth}}{
        \cellcolor[rgb]{0.9,0.9,0.9}{
            \makecell[{{p{\linewidth}}}]{
                \texttt{\tiny{[GM$|$GM]}}
                [1, 1, 1, 0]
            }
        }
    }
    & & \\ \\

    \theutterance \stepcounter{utterance}  
    & & \multicolumn{4}{p{0.6\linewidth}}{
        \cellcolor[rgb]{0.9,0.9,0.9}{
            \makecell[{{p{\linewidth}}}]{
                \texttt{\tiny{[GM$\rangle$P2]}}
                \texttt{You are playing a cooperative negotiation game in which you have to agree on how to divide a set of items among you and another player.} \\
\\ 
\texttt{Rules:} \\
\texttt{(a) You and the other player are given a set of items. Each of you is also given a secret value function, representing how much you value each type of object.} \\
\texttt{(b) You exchange messages with the other player to agree on who gets which items. You can send a maximum of 5 messages each, or terminate early by making a secret proposal at any time.} \\
\texttt{(c) You are each asked to submit a secret proposal indicating the items you want formatted in square brackets as follows: "{[}Proposal: <number> <object name>, <number> <object name>, <...>{]}"} \\
\texttt{(d) If your proposals are complementary, i.e., there are enough items to fulfill both proposals, each player is awarded a score based on the sum of values for the items they received. Otherwise, both of you get zero points.} \\
\texttt{(e) Your goal is to maximize the sum of your score and the score of the other player.} \\
\\ 
\texttt{Let us start.} \\
\\ 
\texttt{The set of available items is:} \\
\\ 
\texttt{2 mirrors, 1 brick, 3 knives, 1 comb} \\
\\ 
\texttt{Your secret value function is:} \\
\\ 
\texttt{mirror: 2, brick: 1, knife: 0, comb: 5} \\
\\ 
\texttt{IMPORTANT: Your messages, unless it is the secret proposal, are directly transmitted to the other player, so do not include any response to the rules or text announcing your message. To make a secret proposal, use the indicated format. Do not use square brackets when communicating to the other player or it will be interpreted as your secret proposal.} \\
\\ 
\\ 
\texttt{The other player has submitted a secret proposal. It is now your turn to submit a proposal of your own.} \\
            }
        }
    }
    & \\ \\

    \theutterance \stepcounter{utterance}  
    & & & \multicolumn{4}{p{0.6\linewidth}}{
        \cellcolor[rgb]{0.9,0.9,0.9}{
            \makecell[{{p{\linewidth}}}]{
                \texttt{\tiny{[GM$\langle$P2]}}
                \texttt{<think>} \\
\texttt{Okay, let's see. I need to figure out how to divide the items between me and the other player to maximize the total score. The available items are 2 mirrors, 1 brick, 3 knives, 1 comb. My value function is mirror: 2, brick: 1, knife: 0, comb: 5. The other player has already submitted a secret proposal, so I need to respond with my own.} \\
\\ 
\texttt{First, I should consider what the other player's proposal might be. Since their proposal is secret, I don't know exactly, but maybe I can infer based on typical strategies. Let's think about possible allocations.} \\
\\ 
\texttt{My highest value is the comb at 5, then mirrors at 2 each. Knives are worth 0 to me, so I probably don't want any. The brick is 1, which is low. So ideally, I want the comb and as many mirrors as possible. But I need to make sure that the other player's proposal doesn't conflict with mine.} \\
\\ 
\texttt{If the other player also values the comb highly, they might have asked for it. But since my value for comb is 5, maybe they have a lower value. Alternatively, maybe they value knives or the brick more. Since I don't know their value function, I have to assume they might have different priorities.} \\
\\ 
\texttt{Let me think about possible allocations. Let's say I take the comb and both mirrors. That would give me 5 + 2 + 2 = 9. Then the remaining items are 1 brick and 3 knives. If the other player's proposal includes the brick and some knives, maybe that's possible. But since knives are 0 to me, I can let them have all the knives. The brick is 1 for me, so maybe they can take that too. So if I take comb, mirrors, and leave the rest, their proposal might be able to take the brick and knives. Let's check the numbers.} \\
\\ 
\texttt{Total items: 2 mirrors, 1 brick, 3 knives, 1 comb. If I take 1 comb, 2 mirrors, that uses up 3 items. The remaining are 1 brick and 3 knives. If the other player's proposal is for 1 brick and 3 knives, then that's possible. But their proposal might be different. However, since I don't know their values, I need to make sure that my proposal} \\
            }
        }
    }
     \\ \\

    \theutterance \stepcounter{utterance}  
    & & & \multicolumn{2}{p{0.3\linewidth}}{
        \cellcolor[rgb]{0.9,0.9,0.9}{
            \makecell[{{p{\linewidth}}}]{
                \texttt{\tiny{[GM$|$GM]}}
                \texttt{continue} \\
            }
        }
    }
    & & \\ \\

    \theutterance \stepcounter{utterance}  
    & & & \multicolumn{2}{p{0.3\linewidth}}{
        \cellcolor[rgb]{0.9,0.9,0.9}{
            \makecell[{{p{\linewidth}}}]{
                \texttt{\tiny{[GM$|$GM]}}
                \texttt{player was instructed to make a proposal} \\
            }
        }
    }
    & & \\ \\

\end{supertabular}
}

\end{document}
